\chapter{Boolean Analysis — Linear}
%%% generated by the blueprint pipeline from TCSlib.BooleanAnalysis.GaussianSpace.Noise.Linear
\section{Overview}

Declarations of the Lean module \texttt{TCSlib.BooleanAnalysis.GaussianSpace.Noise.Linear}.

\section{Declarations}

\begin{lemma}[Measure preserving gaussian neg]
\lean{BooleanAnalysis.GaussianSpace.measurePreserving_gaussianNeg}
\statementsource{vershynin-hdp-book}{
  pdf-681b6f3d947f-p058-b003,
  pdf-681b6f3d947f-p058-b004
}
\label{BooleanAnalysis.GaussianSpace.measurePreserving_gaussianNeg}
\leanok
\uses{BooleanAnalysis.ThresholdFunctions.standardGaussianMeasure}
\difficulty{3}
% dataset-pipeline: stub -- prose pending build_informal_statements.py
Statement pending regeneration.
\end{lemma}

\begin{lemma}[Gaussian invariance under paired reflection]
\lean{BooleanAnalysis.GaussianSpace.measurePreserving_gaussianReflection}
\label{BooleanAnalysis.GaussianSpace.measurePreserving_gaussianReflection}
\leanok
\uses{BooleanAnalysis.GaussianSpace.measurePreserving_gaussianNeg, BooleanAnalysis.GaussianSpace.measurePreserving_gaussianRotation, BooleanAnalysis.ThresholdFunctions.standardGaussianMeasure}
\difficulty{4}
Let $n$ be a natural number and let $\theta\in\mathbb R$. For independent vectors
$z,w\in\mathbb R^n$ with the product standard Gaussian law, the map sending $(z,w)$ to
the pair whose coordinates are $((\cos\theta)z_i+(\sin\theta)w_i)_{i=1}^n$ and
$((\sin\theta)z_i-(\cos\theta)w_i)_{i=1}^n$ preserves that product law.
\end{lemma}

\begin{lemma}[Gaussian characteristic functional]
\lean{BooleanAnalysis.GaussianSpace.charFunDual_standardGaussian}
\label{BooleanAnalysis.GaussianSpace.charFunDual_standardGaussian}
\leanok
\uses{BooleanAnalysis.ThresholdFunctions.standardGaussianMeasure}
\difficulty{4}
Let $\gamma$ be the standard Gaussian probability measure on $\mathbb R$, and let
$L:\mathbb R\to\mathbb R$ be a continuous real-linear functional. Its characteristic
functional at $L$ is $\exp\!\left(-((L(1))^2)/2\right)$, viewed as a complex number.
\end{lemma}

\begin{lemma}[Characteristic function of a Gaussian linear functional]
\lean{BooleanAnalysis.GaussianSpace.charFunDual_gaussianPi}
\label{BooleanAnalysis.GaussianSpace.charFunDual_gaussianPi}
\leanok
\uses{BooleanAnalysis.GaussianSpace.charFunDual_standardGaussian, BooleanAnalysis.ThresholdFunctions.standardGaussianMeasure}
\difficulty{3}
Let $n$ be a natural number, and let $L:\mathbb R^n\to\mathbb R$ be a continuous linear
functional. For a vector $Z$ with independent standard Gaussian coordinates, the
characteristic function of $L(Z)$ is
\[
\mathbb E\bigl[e^{iL(Z)}\bigr]=\exp\!\left(-\frac12\sum_{j=1}^n L(e_j)^2\right),
\]
where $e_j$ are the standard coordinate vectors.
\end{lemma}

\begin{lemma}[Expansion of squared projection coefficients]
\lean{BooleanAnalysis.GaussianSpace.sum_sq_projectionCoefficients}
\label{BooleanAnalysis.GaussianSpace.sum_sq_projectionCoefficients}
\leanok
\difficulty{3}
Let $n\ge 0$, let $a,b:\{0,\ldots,n-1\}\to\mathbb R$, and let $u,v\in\mathbb R$. Then
\[\sum_{i=0}^{n-1}(u a_i+v b_i)^2=u^2\sum_{i=0}^{n-1}a_i^2+2uv\sum_{i=0}^{n-1}a_i b_i+v^2\sum_{i=0}^{n-1}b_i^2.\]
\end{lemma}

\begin{lemma}[Linear functional on a pair]
\lean{BooleanAnalysis.GaussianSpace.realLinearFunctional_pair}
\label{BooleanAnalysis.GaussianSpace.realLinearFunctional_pair}
\leanok
\difficulty{4}
Let $L:\bbr^2\to\bbr$ be a continuous real-linear functional. For all $x,y\in\bbr$,
\[
L(x,y)=xL(1,0)+yL(0,1).
\]
\end{lemma}

\begin{lemma}[Characteristic function of a Gaussian vector]
\lean{BooleanAnalysis.GaussianSpace.charFun_gaussianEuclidean}
\label{BooleanAnalysis.GaussianSpace.charFun_gaussianEuclidean}
\leanok
\uses{BooleanAnalysis.ThresholdFunctions.standardGaussianMeasure}
\difficulty{3}
Let $n$ be a natural number, and let $X$ be the random vector in $\mathbb R^n$ whose
coordinates are independent standard Gaussian random variables. For every $t\in\mathbb
R^n$, the characteristic function of $X$ at $t$ is $\exp(-\|t\|^2/2)$.
\end{lemma}

\begin{lemma}[Gaussian measure isometry invariance]
\lean{BooleanAnalysis.GaussianSpace.measurePreserving_gaussianEuclideanIsometry}
\proofsource{vershynin-hdp-book}{pdf-681b6f3d947f-p058-b004}
\label{BooleanAnalysis.GaussianSpace.measurePreserving_gaussianEuclideanIsometry}
\leanok
\uses{BooleanAnalysis.GaussianSpace.charFun_gaussianEuclidean, BooleanAnalysis.ThresholdFunctions.standardGaussianMeasure}
\proofstep
  {BooleanAnalysis.ThresholdFunctions.standardGaussianMeasure}
  {context}
  {vershynin-hdp-book}
  {pdf-681b6f3d947f-p058-b002}
\proofstep
  {BooleanAnalysis.GaussianSpace.charFun_gaussianEuclidean}
  {context}
  {vershynin-hdp-book}
  {pdf-681b6f3d947f-p058-b002,
    pdf-681b6f3d947f-p058-b005}
\difficulty{4}
Let $n$ be a natural number and let $Q$ be a linear isometry from Euclidean $n$-space to
itself. The pushforward of standard Gaussian product measure, viewed as a measure on
Euclidean $n$-space, under $Q$ is that same measure.
\end{lemma}

\begin{lemma}[Measure preserving gaussian isometry]
\lean{BooleanAnalysis.GaussianSpace.measurePreserving_gaussianIsometry}
\statementsource{vershynin-hdp-book}{
  pdf-681b6f3d947f-p058-b003,
  pdf-681b6f3d947f-p058-b004
}
\label{BooleanAnalysis.GaussianSpace.measurePreserving_gaussianIsometry}
\leanok
\uses{BooleanAnalysis.GaussianSpace.measurePreserving_gaussianEuclideanIsometry, BooleanAnalysis.ThresholdFunctions.standardGaussianMeasure}
\difficulty{3}
% dataset-pipeline: stub -- prose pending build_informal_statements.py
Statement pending regeneration.
\end{lemma}

\begin{lemma}[Gaussian law of a correlated endpoint]
\lean{BooleanAnalysis.GaussianSpace.measurePreserving_gaussianEndpoint}
\proofsource{vershynin-hdp-book}{
  pdf-681b6f3d947f-p058-b003,
  pdf-681b6f3d947f-p058-b004
}
\label{BooleanAnalysis.GaussianSpace.measurePreserving_gaussianEndpoint}
\leanok
\uses{BooleanAnalysis.GaussianSpace.measurePreserving_gaussianRotation, BooleanAnalysis.ThresholdFunctions.standardGaussianMeasure}
\proofstep
  {BooleanAnalysis.ThresholdFunctions.standardGaussianMeasure}
  {context}
  {vershynin-hdp-book}
  {pdf-681b6f3d947f-p058-b002,
    pdf-681b6f3d947f-p119-b003}
\proofstep
  {BooleanAnalysis.GaussianSpace.measurePreserving_gaussianRotation}
  {direct}
  {vershynin-hdp-book}
  {pdf-681b6f3d947f-p058-b003,
    pdf-681b6f3d947f-p058-b004}
\difficulty{1}
Let $n$ be a natural number and let $\rho\in[-1,1]$. If $X,Y$ are independent standard
Gaussian vectors in $\mathbb R^n$, then $\rho X+\sqrt{1-\rho^2}\,Y$ has the standard
Gaussian product law on $\mathbb R^n$.
\end{lemma}

\begin{lemma}[Gaussian law of a linear combination]
\lean{BooleanAnalysis.GaussianSpace.measurePreserving_gaussianCombination}
\proofsource{vershynin-hdp-book}{
  pdf-681b6f3d947f-p058-b004,
  pdf-681b6f3d947f-p170-b003
}
\label{BooleanAnalysis.GaussianSpace.measurePreserving_gaussianCombination}
\leanok
\uses{BooleanAnalysis.GaussianSpace.measurePreserving_gaussianEndpoint, BooleanAnalysis.ThresholdFunctions.standardGaussianMeasure}
\proofstep
  {BooleanAnalysis.ThresholdFunctions.standardGaussianMeasure}
  {context}
  {vershynin-hdp-book}
  {pdf-681b6f3d947f-p058-b002}
\proofstep
  {BooleanAnalysis.GaussianSpace.measurePreserving_gaussianRotation}
  {direct}
  {vershynin-hdp-book}
  {pdf-681b6f3d947f-p058-b004}
\proofstep
  {BooleanAnalysis.GaussianSpace.measurePreserving_gaussianEndpoint}
  {context}
  {vershynin-hdp-book}
  {pdf-681b6f3d947f-p058-b004}
\difficulty{4}
Let $n$ be a natural number, and let $b,c\in\mathbb R$ satisfy $0\le c$ and $b^2+c^2=1$.
If $X,Y$ are independent standard Gaussian vectors in $\mathbb R^n$, then $bX+cY$ has
the standard Gaussian product law on $\mathbb R^n$.
\end{lemma}

\begin{lemma}[Products preserve the correlation interval]
\lean{BooleanAnalysis.GaussianSpace.correlation_mul_mem}
\label{BooleanAnalysis.GaussianSpace.correlation_mul_mem}
\leanok
\difficulty{3}
If $\rho,\sigma\in[-1,1]$, then $\rho\sigma\in[-1,1]$.
\end{lemma}

\begin{lemma}[Gaussian noise increment variance identity]
\lean{BooleanAnalysis.GaussianSpace.gaussianNoise_increment_variance}
\label{BooleanAnalysis.GaussianSpace.gaussianNoise_increment_variance}
\leanok
\uses{BooleanAnalysis.GaussianSpace.correlation_mul_mem}
\difficulty{3}
Let $n$ be a natural number and let $\rho,\sigma\in[-1,1]$. Then
\[\bigl(\sigma\sqrt{1-\rho^2}\bigr)^2+\bigl(\sqrt{1-\sigma^2}\bigr)^2=\bigl(\sqrt{1-(\rho\sigma)^2}\bigr)^2.\]
\end{lemma}

\begin{lemma}[Zero correlation gives independent Gaussian vectors]
\lean{BooleanAnalysis.GaussianSpace.correlatedGaussianMeasure_zero}
\proofsource{boolean-ch11-gaussian-invariance}{pdf-8b2c50cf9621-p004-b001}
\label{BooleanAnalysis.GaussianSpace.correlatedGaussianMeasure_zero}
\leanok
\uses{BooleanAnalysis.GaussianSpace.correlatedGaussianMeasure, BooleanAnalysis.ThresholdFunctions.standardGaussianMeasure}
\proofstep
  {BooleanAnalysis.ThresholdFunctions.standardGaussianMeasure}
  {direct}
  {boolean-ch11-gaussian-invariance}
  {pdf-8b2c50cf9621-p002-b001}
\proofstep
  {BooleanAnalysis.GaussianSpace.correlatedGaussianMeasure}
  {direct}
  {boolean-ch11-gaussian-invariance}
  {pdf-8b2c50cf9621-p004-b001}
\difficulty{1}
For every nonnegative integer $n$, the correlated Gaussian measure with correlation
parameter $0$ equals the product of two standard Gaussian product measures on $\mathbb
R^n$.
\end{lemma}
